\section{Coupled shared cores and the remaining power thresholds}

\subsection*{Finding M-016: even jets from a matrix of root bridges}
Let \(r\geq4\), \(\ell\geq1\), and put
\[
 R=r-1,\qquad b=2\ell+2,\qquad h=\lceil\sqrt r\rceil,
 \qquad
 L=2\binom{r+\ell-1}{\ell}-r,\qquad N=hL.
\]
For every \(n\geq N\), the determinant map from
\(n\)-by-\(n\) linear pencils in \(r\) variables dominates
the whole coefficient jet through degree \(b\).
Its differential admits an explicit integral right inverse.
In particular, in characteristic zero,
\begin{equation}\label{eq:M016K}
 \lambda\vdash nd,\quad\ell(\lambda)\leq r,\quad\lambda_2\leq b
       \quad\Longrightarrow\quad K_n(\lambda)=a_\lambda(d[n]).
\end{equation}
The useful small orders are
\[
       b=4:\ N=\lceil\sqrt r\rceil r,
       \qquad b=6:\ N=\lceil\sqrt r\rceil r^2.
                                                               \tag{M-016.1}
\]
The construction concerns detected coefficients. It asserts
neither exact representation of each truncated form nor
an algebraic section of the full determinant map.

\paragraph{The components and their root matrix.}
Use the shared-core graph of F-022 with variables
\(x,z_1,\ldots,z_R\) and cyclic permutation
\(\sigma=(1\ 2\ \cdots\ R)\). It has one root,
\(R\) shared degree-one vertices, and separate inward and
outward monomial-tree vertices of degrees two through \(\ell\).
Its size is \(L\). For \(\ell=1\), only the root and
shared degree-one vertices occur.
Take \(h\) disjoint copies, labeled \(1,\ldots,h\),
with roots \(o_1,\ldots,o_h\).

Since \(R\leq h^2-1\), inject the \(R\) variable labels
into the \(h^2-1\) ordered pairs \((a,c)\ne(h,h)\).
Let \(C_{ac}\) be the assigned variable, or zero when
that pair is unused. Add a directed bridge
\(o_a\longrightarrow o_c\) labeled \(C_{ac}\) for each
nonzero label. Diagonal bridges are allowed: they add a
linear term to that root's diagonal matrix entry.
In particular \(C_{hh}=0\).

Let \(B(z)\) be the resulting adjacency matrix and
\(A_0=xI-B(z)\). At \(x=1\), eliminate the nonroot
blocks, whose directed paths outside the roots have no cycles.
The shared degree-one star in each component contributes
the same quadratic polynomial
\[
                 s(z)=\sum_{i=1}^R z_i z_{\sigma(i)}.
\]
Thus the Schur complement on the roots is exactly
\[
                         W=I_h-C-sI_h,
                 \qquad \det A_0(1,z)=\det W.           \tag{M-016.2}
\]
All eliminated diagonal entries equal one; the higher tree
vertices cannot contribute a closed path. These observations
also prove (M-016.2) directly by the block determinant formula.

\paragraph{Cofactors detected by the root matrix.}
The pure monomial cofactors in F-022 can be chosen to use
the central root, with no path bypassing it. Consequently,
for every monomial \(z^\gamma\) of degree at most \(2\ell\)
and every component \(a\), an elementary coefficient
direction has cofactor
\[
                         z^\gamma\operatorname{adj}(W)_{aa}.
                                                               \tag{M-016.3}
\]
For degree zero this is the root diagonal cofactor itself.
For degree at most \(\ell\), use the path from the root
to an outward monomial vertex. For larger degrees use
F-022's split into two monomials, each of degree at most
\(\ell\), with distinct shared parent vertices. That
condition removes the direct path bypassing the root.
The local path factors multiply to \(z^\gamma\);
the remaining inverse entry is \((W^{-1})_{aa}\).
Multiplying by \(\det W\) gives (M-016.3).

For distinct components \(a,c\), and monomials
\(z^\alpha,z^\beta\), each of degree exactly \(\ell\),
the analogous cofactor from the inward vertex in component
\(a\) to the outward vertex in component \(c\) is
\[
                     z^{\alpha+\beta}\operatorname{adj}(W)_{ac}.
                                                               \tag{M-016.4}
\]
Different components have no nonroot vertices in common,
so there is no bypass term. With the adjacency convention
that \(B_{uv}\) labels \(u\to v\), the derivative direction
is the reverse matrix entry: its cofactor is
\(\det A_0(A_0^{-1})_{uv}\). This fixes the orientation
in (M-016.4).

The inverse calculations may first be made over
\(\mathbb Q(z)\), where the constant term one makes
\(W\) invertible. After multiplying by the determinant
they are identities of integral polynomials, and do not
restrict the coefficient parameter points.

\paragraph{First terms of the root cofactors.}
Let \(O(z^j)\) mean terms of total degree at least \(j\).
Since \(C\) is linear and \(s\) is quadratic, expansion
of the determinant and inverse at zero gives
\[
 \operatorname{adj}(W)
     =I_h+C-\operatorname{tr}(C)I_h+O(z^2).
\]
It follows that
\[
 \begin{split}
 \operatorname{adj}(W)_{ac}&=C_{ac}+O(z^2)\quad(a\ne c),\\
 \operatorname{adj}(W)_{aa}-\operatorname{adj}(W)_{hh}
                 &=C_{aa}+O(z^2),\\
 \operatorname{adj}(W)_{hh}&=1+O(z).
 \end{split}                                             \tag{M-016.5}
\]
These first terms can also be read directly from the
cofactor permutation expansion; their coefficients are integral.

\begin{proof}[Proof of the surjective integral differential]
Consider a target monomial of degree \(b=2\ell+2\).
Choose one of its variables as a bridge label and another
as the linear factor \(H\) in an entry variation.
The remaining multiset has degree \(2\ell\).

If the chosen bridge variable is in an off-diagonal slot
\((a,c)\), split the remaining multiset into \(\alpha,\beta\)
of degree \(\ell\). By (M-016.4) and (M-016.5),
the entry variation by \(H\) has, after truncation through
degree \(b\), exactly the chosen target monomial.
Every other term has degree at least \(2\ell+3\).

If the bridge variable is in a diagonal slot \((a,a)\),
it has \(a\ne h\). Write the remaining monomial as
\(z^\gamma\), \(|\gamma|=2\ell\). Take the difference
of the two elementary directions in (M-016.3) for components
\(a,h\), both varied by \(H\). The second line of
(M-016.5) again gives exactly the target monomial through
degree \(b\). Thus the differential contains the entire
top homogeneous coefficient space, over \(\mathbb Z\).

For degrees at most \(2\ell\), (M-016.3) in component
\(h\), varied by \(H=x\), has lowest term precisely
\(z^\gamma\) with coefficient one. Degree \(2\ell+1\)
uses a degree-\(2\ell\) cofactor times a variable.
Starting at degree \(b-1\) and proceeding downwards,
subtract integral combinations of already obtained higher
monomials. This yields every homogeneous target monomial
exactly. No division occurs.

We have therefore constructed integral linear combinations
of source coefficient directions which map to the monomial
basis of the full jet module. They give a right inverse to
the differential over \(\mathbb Z\). Restrict the pencil
parameters to the affine space generated by those directions
through \(A_0\), and translate its target by the coefficient
vector of \(\det W\). The restricted polynomial map is
\(F(t)=t+O(t^2)\). The lowest homogeneous term argument
of M-013 proves injectivity of its polynomial pullback after
every commutative base change. In particular it proves
dominance in characteristic zero.

For \(n>N\), append \(xI_{n-N}\); after \(x=1\)
the same directions and differential apply. Finally F-016
puts every highest-weight function with \(\lambda_2\leq b\)
in the detected coefficient algebra. Injective pullback
annihilates no nonzero determinant-ideal vector there.
This proves \eqref{eq:M016K} for every multiplicity copy.
\end{proof}

\paragraph{A complete second-row gate at integral power thresholds.}
For every pair of integers \(m\geq2,c\geq1\), an actual
internally padded permanent obstruction at \(n\geq m^c\)
requires
\begin{equation}\label{eq:M016powers}
                              \lambda_2\geq c+2.
\end{equation}
For even \(c\geq2\) this is F-022. For odd
\(c=2\ell+1\geq3\), a positive permanent multiplicity
has length \(r\leq m^2\). The cases \(r\leq3\)
are already excluded by O-006--O-007. For \(r\geq4\),
the size above is at most
\[
 m\left(2\binom{m^2+\ell-1}{\ell}-m^2\right).
\]
The monotonicity used here follows from
\(L(r+1)-L(r)=2\binom{r+\ell-1}{\ell-1}-1\geq1\).
At \(\ell=1\) and \(\ell=2\) this is exactly
\(m^3\) and \(m^5\), respectively. At \(\ell\geq3\),
M-013.2 bounds it by \((5/8)m^{2\ell+1}\).
Thus \(n\geq m^c\) gives full determinant multiplicity
whenever \(\lambda_2\leq c+1\), proving the assertion.

For \(c=1\), in fact every ambient type with second row
at most two has full determinant multiplicity. F-014 and
the quadratic multiplicity-free decomposition in O-013
show that its only nonzero non-one-row types are
\((nd-2(r-1),2^{r-1})\), each with ambient multiplicity one.
The nonzero determinant Hessian section in O-013 realizes
that copy for every allowed \(r\leq\min(d,n^2)\).
One-row types also have full multiplicity. Hence no actual
obstruction has \(\lambda_2\leq2\), completing
\eqref{eq:M016powers}.

Equivalently, every actual obstruction with second row
\(b=\lambda_2\) satisfies the useful degree-independent
restriction \(n<m^{b-1}\). It is a necessary restriction,
not a sufficient test for an obstruction. At \(c=25\)
it recovers \(\lambda_2\geq27\); the new endpoint
ranges here are \(c=1,2,3,4,5\).

\paragraph{Grothendieck dictionary.}
The root Schur complement and cofactor formulas are
identities over \(\mathbb Z\). The finite jet module is
locally free, and the explicit integral splitting of the
differential produces a universally injective polynomial
coordinate pullback. Its monomial supports are torus
comodules under SGA~3, I, Proposition~4.7.3, pp.~35--36.
The Schur-copy support used in the final characteristic-zero
step is SGA~6, VI, Theorem~3.3, pp.~392--393, as implemented
by F-016's Cauchy flag-minor argument. The resulting
vanishing concerns the entire ideal comodule, in the sense
of SGA~3, I, Proposition~4.7.2, pp.~34--35.
Sharing and coupling vertices retains explicit coordinate
maps; an equality of virtual \(K_0\) classes would not
provide the differential or its integral splitting.
